% Folder 77, batch 05 — archivists' pages 81–100.
% Pass: Opus 5.5 (claude-opus-5-5), 2026-10-03 — first pass, unchecked against the pages by a human.
%
% Page 82 is an unrelated typed verso (administrative notice) and is skipped.
% Page 87 was scanned upside down; it is read turned.
\documentclass[11pt,a4paper]{article}
\input{../preamble/grothendieck}

\folder{77}
\batch{5}
\pages{81}{100}
\foldertitle{Quadriques affines, extensions quadratiques : notes manuscrites (s.d.).}
\dating{[à partir de 1977]}
\watermark{Édition de démonstration}

\begin{document}

\section{Extensions quadratiques étales et discriminant}

\page{81}
\note{la page continue un argument commencé avant le lot : la première ligne reprend au milieu d'une phrase.}

\ill{} $A_\Delta[T]/(T^2+bT+c)$ est une extension quadratique étale de $A_\Delta$, et sa restriction à $A_{2\Delta}$ est isomorphe à $A_{2\Delta}[U]/(U^2-\Delta)$, en posant
\[
U=\frac{1}{2u}\Bigl(T+\frac{b}{2}\Bigr).
\]

\struck{Dém}

\struck{Lemme} \struck{Supposons $A_\Delta$ \ill{} \ill{}}

\marginal{\uncertain{dès que} \ill{} \uncertain{connexes}}
Pour que $\exists\, B_\Delta$ ext.\ quadratique étale de $A_\Delta$, se restreignant sur $A_{2\Delta}$ suivant $A_{2\Delta}[U]/(U^2-\Delta)$, il faut et il suffit que l'on puisse trouver $b,c\in A_\Delta$ et $\nu\in\mathbf{N}$ et $u\in A_\Delta^{*}$ tels que
\[
b^2-4c=\Delta^{\nu}u^2 .
\]
\note{un exposant écrit sur $\Delta$ est biffé, remplacé par $\nu$.}

Or, quitte à multiplier $b$, $c$ par des \add{carrés} \uncertain{de} $\Delta$ ($\Delta^{\alpha}$ et $\Delta^{2\alpha}$ resp.), puis \ill{}, on peut supposer $b,c\in A$, \struck{\ill{}} $u\in A$.

\struck{Supposons que}
\[
\text{\struck{$A_\Delta^{*}=\{u\Delta^{d}\mid d\in\mathbf{Z},\ u\in A^{*}\}$}}
\]
\note{ce cadre et ce qui suit jusqu'au bas de la page sont barrés de traits obliques.}
\struck{donc on aurait}
\[
\text{\struck{$b^2-4c=\Delta^{2d+1}u^2\quad (d\in\mathbf{Z},\ u\in A^{*})$}}
\]
\struck{On a $d\in\mathbf{N}$, \ill{} \ill{} $d\geqslant 0$, \ill{}, on aurait $b^2-4c=\Delta^{\nu}u^2$ $(\nu\geqslant 0)$, $\Delta^{-\nu}(b^2-4c)=u^2$}

\page{83}
\drawing{en tête de page, un schéma de traits obliques et de courbes, portant les inscriptions « dans $\Delta^{\nu}$ \struck{$A^{*}$} », « $\partial\in A^{*}$ », « $\Delta\notin A^{*}$ », « \ill{} » ; l'ensemble est barré.}

Nous voulons prouver que $u_0$ est \uncertain{inversible}. \struck{\ill{} $A_0$ \ill{} \ill{}} \uncertain{Il suffit} \ill{} \uncertain{régulier} dans $A_0$ ; \ill{} qui ne \ill{} en aucune composante \ill{} réduite de $A_0$, ce qui résulte de \ill{} ce qui précède.

\drawing{dans la marge gauche, deux droites horizontales marquées $\Delta=0$ coupées par une courbe marquée $2=0$.}

D\ill{} \ill{} \ill{}

\textbf{Proposition}. $A_0=A/2A$. On suppose \add{$2$ $A$-régulier,} $A$ \ill{} \ill{}, $\Delta\in A$. On suppose \add{$A_0$ réduit et} \ill{} \ill{} des fractions, \ill{} \ill{} singulier de $A_0$, \struck{$A_0$}

$\Delta_0$ ($=\Delta$ mod $2$) \ill{} \ill{} \ill{} \ill{}

\[
\begin{cases}
\operatorname{Ker}\bigl(\operatorname{Pic}(A_\Delta)\to\operatorname{Pic}(A_{2\Delta})\bigr)=0, \\
A_\Delta \text{ réduit et } \dots
\end{cases}
\]
\note{un symbole biffé et illisible dans cette formule est omis.}
\note{la seconde ligne de l'accolade porte en fin la mention « les \ill{} de $V(A_0)$ ».}

\marginal{il suffit que $A$ normal et $A_0$ intègre}
\uncertain{Conditions} \struck{des} \struck{$A_{2\Delta}$} \ill{} \uncertain{équivalentes} :

\struck{Pour que} il existe $b,c\in A$ avec
\[
b^2-4c=\Delta
\]

\struck{il faut} b) il existe une extension quadratique étale sur $A_\Delta$ de $A$, \uncertain{qui a réduite} sur $A_{2\Delta}$ \uncertain{donne} l'extension étale
\[
A_{2\Delta}[U]/(U^2-\Delta)\qquad [\text{NB. } A[T]/(T^2+bT+c)]
\]

b$'$) il existe une extension \add{finie} étale de $A_\Delta$ qui induise sur \struck{\ill{}} $A_{2\Delta}$ \add{\ill{}} \ill{} $A_{2\Delta}[U]/(U^2-\Delta)$,

\marginal{NB il suffit étale aux pts maximaux de \struck{\ill{}} $\operatorname{Spec}A_0$}

\page{84}
Sous ces conditions, dans b), \ill{} \ill{} nécessaire que l'extension soit étale en les pts de $A_\Delta$

\page{85}
\[
\text{\struck{$\begin{cases} 2\in A \text{ régulier} \\ A_0=A/2A \text{ intègre normal} \\ A \text{ noeth.} \\ \Delta\in A,\quad \Delta\notin 2A \end{cases}$}}
\]
\note{au-dessus de « régulier », un mot ajouté est illisible.}

\textbf{Proposition}. Soit $A$ anneau noeth., $\Delta\in A$\struck{$-2A$}. On suppose que $2$ est $A$-régulier, que $A_0=A/2A$ est intègre et normal, que $\Delta\notin 2A$. Conditions équivalentes :

a) $\exists\, b,c\in A$ avec
\[
b^2-4c=\Delta
\]

b) $\exists$ extension quadratique $B$ de $A$ telle que $B_\Delta$ soit étale sur $A_\Delta$, et telle que $B_{2\Delta}\simeq A_{2\Delta}[U]/(U^2-\Delta)$.

c) $\exists$ extension [quadratique] \add{finie} \struck{étale} $B_\Delta$ de $A_\Delta$, telle que $(B_\Delta)_2\simeq A_{2\Delta}[U]/(U^2-\Delta)$.

De plus, l'extension \struck{($B$ de b)} \add{$B_\Delta$ de c)} \struck{\ill{}} est déterminée à isom.\ unique près.

\page{86}
$L$.

\[
\begin{cases}
X\supset U=X-D \\
X_0=V(2\mathcal{O}_X)\qquad X_0 \text{ intègre normal}
\end{cases}
\]
\[
\ U-U_0=U\cap V\qquad V=X-X_0,\qquad U_0=U\cap X_0 .
\]
\note{un symbole biffé et illisible dans cette formule est omis.}

$U'$ ext.\ quadratique de $U$ :
\[
L_{U'}^{\,\otimes 2}\xrightarrow{\ \sim\ }\underline{\mathcal{O}}_{U'}
\]
\note{un symbole biffé et illisible dans cette formule est omis.}
$L_{U'}$ se prolonge à $L_U$.
\[
\begin{cases}
E_U \text{ fibré sur } L_U \\
L_U^{\otimes 2}\ \text{\struck{$\simeq\mathcal{O}_S$}}\ \Delta
\end{cases}
\]
$L_U$ se prolonge en $L$ \emph{donné}, \struck{$L^{\otimes2}\simeq\mathcal{O}_S$}, $\Delta$ section de $L^{\otimes 2}$.
\note{deux symboles griffonnés et biffés suivent à droite.}

$V=X\smallsetminus X_0$
\[
U\cup V=W=X\smallsetminus \underbrace{X_0\cap D}_{D_0}
\]
\[
\widetilde{V}_1\to V_1\subset L
\]
\emph{OK} sur $V=X-{}$\struck{$D$}${}X_0$.
\emph{Mais} pas sur les \uncertain{pts} de $X_0$.

En \ill{} \ill{} $X_0\cap U=U-U\cap V$, on \ill{} pour les \ill{} \ill{} \ill{} \ill{}

\[
\tfrac{1}{2}L\qquad\qquad
\begin{cases} 2x\in L \\ x^2 \end{cases}
\]
\note{la page s'arrête sur cette accolade inachevée.}

\page{87}
\note{calcul d'un développement en termes de graphes à quatre sommets ; la page entière est barrée d'un long trait oblique. Les colonnes donnent, pour chaque type de monôme, un dessin du graphe et deux décomptes.}

\[
\begin{array}{lccc}
2^4\,Y_1Y_2Y_3Y_4 & \text{(quatre points)} & 1 & 1 \\
-2^2\,Y_1Y_2(Y_{3,4})^2 & \text{(deux points, une arête)} & 6 & 6 \\
2^2\,Y_1\,Y_{23}Y_{34}Y_{41} & \text{(un point, un triangle)} & 4 & 8 \\
(Y_{12}Y_{34})^2 & \text{(deux arêtes disjointes)} & 3 & 3 \\
-2\,Y_{12}Y_{23}Y_{34}Y_{41} & \text{(un carré)} & 3 & 6 \\
 & & & 24=4!
\end{array}
\]
\note{le coefficient $2^4$ de la première ligne est écrit sur un autre chiffre ; dans la dernière colonne, le premier « 1 » est écrit sur un chiffre biffé.}

\[
\Bigl(\sum Y_{12}Y_{34}\Bigr)^2=\sum
\]
\note{l'indice de cette somme est illisible.}
\[
(-1)\sum_{\gamma\in\Gamma_0}M_\gamma
\ +\ (-1)^{\nu_2(\gamma)+1}\,2\!\!\sum_{\substack{\gamma\in\Gamma\\ \nu^{*}(\gamma)=1\\ \nu_1(\gamma)=0}}\!\! M_\gamma
\]
\note{le signe « $+$ » central est surchargé ; le facteur $2$ est ajouté devant le second $\Sigma$.}

\drawing{deux rangées de tirets, et un hexagone.}

\[
(-1)^n\Delta=\sum_{\gamma\in\Gamma_0}M_\gamma+(-1)^{n+\nu_2(\gamma)}\sum_{\gamma\in\Gamma_1}M_\gamma
\]
\note{un trait relie l'exposant $\nu_2(\gamma)$ à l'indice de sommation $\gamma\in\Gamma_1$.}
\[
(-1)^n\Delta-B^2\equiv 2\sum_{\gamma\in\Gamma_1}(-1)^{n+1+\nu_2(\gamma)}M_\gamma
\ -\underbrace{\sum_{\substack{\gamma,\gamma'\in\Gamma_0\\ \gamma\neq\gamma'}}M_\gamma M_{\gamma'}}_{2\sum_{\gamma\in\Gamma_1}M_\gamma}
\]
\[
=2\sum_{\gamma\in\Gamma_1}\bigl[(-1)^{(\dots)}-1\bigr]M_\gamma
\]
\note{l'exposant de $(-1)$ est illisible ; sous le crochet, une flèche vers l'exposant et la mention « \uncertain{pair} », puis « $0$ \ill{} $2$ ».}

\section{Quadriques affines}

\page{88}
Quadriques affines \struck{et invariants des polyèdres réguliers sphériques\,…}
\note{titre écrit sur une chemise de papier rose ; la seconde moitié est biffée.}

\page{89}
« Invariants » des groupes des polyèdres \emph{réguliers} sphériques

\textbf{1}. Soit $S$ schéma de base, $(E,q,\omega)$ un Module spécial quadratique sur $S$ de rang $3$, donc
\[
\delta(q)=-\omega^{\otimes 2}\qquad \bigl(\omega\in\Gamma(S,(\det E)^{*})\bigr) \tag{1}
\]
On sait alors que $E$ est muni d'une loi $[\ ,\ ]$, plus précisément, on peut trouver une algèbre de quaternions tordue $\mathcal{M}$ sur $S$ (définie à isom.\ unique près) et un isomorphisme de $\mathcal{O}_S$-Mod.
\[
E\xleftarrow{\ \sim\ }\gamma\mathcal{M}=\operatorname{Ker}\bigl(\mathcal{M}\xrightarrow{\ \mathrm{tr}\ }\underline{\mathcal{O}}_S\bigr) \tag{2}
\]
\note{sous la flèche $\mathrm{tr}$ : « trace \emph{réduite} ».}
transformant $q$ en la forme
\[
q_{\mathcal{M}}(u)=-\det u\qquad \bigl(u\in\Gamma(-,\gamma\mathcal{M})\bigr) \tag{3}
\]
et $\omega$ en la forme $\omega_{\mathcal{M}}$ définie par la condition que
\[
\varphi_{\mathcal{M}}([u,v],w)\,\omega_{\mathcal{M}}=2\,u\wedge v\wedge w\qquad \bigl(u,v,w\in\Gamma(-,\gamma\mathcal{M})\bigr) \tag{4}
\]
\marginal{$\omega$ défini \uncertain{à $\pm$ près} par (4) \ill{} \ill{} que $\varphi_{\mathcal{M}}$ \ill{}}
où $\varphi_{\mathcal{M}}$ est la forme \struck{\ill{}} \add{bilinéaire} sym.\ associée à $q_{\mathcal{M}}$
\[
\begin{cases}
q_{\mathcal{M}}(u+v)=\text{\struck{$\det$}}\ q_{\mathcal{M}}(u)+q_{\mathcal{M}}(v)+\varphi_{\mathcal{M}}(u,v) \\
\text{i.e.}\quad \det(u+v)=\det(u)+\det(v)-\varphi_{\mathcal{M}}(u,v)
\end{cases} \tag{5}
\]
pour $u,v\in\Gamma(-,\gamma\mathcal{M})$, ce qui donne aussi
\[
\varphi_{\mathcal{M}}(u,v)=\operatorname{Tr}uv\qquad u,v\in\Gamma(-,\gamma\mathcal{M}) \tag{6}
\]

\page{90}
de sorte que (4) s'écrit aussi
\[
\text{\struck{$\operatorname{Tr}([u,v]w)\,\omega$}}\quad 2\,u\wedge v\wedge w=\operatorname{Tr}([u,v]w)\,\omega_{\mathcal{M}} \tag{7}
\]
\[
\Bigl(=\operatorname{Tr}(uvw-vuw)\,\omega_{\mathcal{M}},\qquad \operatorname{Tr}(vuw)=-\operatorname{Tr}uvw\Bigr)
\]
\marginal{NB car \ill{} \ill{} \struck{\ill{}} $\omega_{\mathcal{M}}$ indépendant \ill{} le cond.\ $2$ inv.\ \ill{}\,…}
\struck{ou encore}, on a encore
\[
u\wedge v\wedge w=\operatorname{Tr}(uvw)\,\omega_{\mathcal{M}}\qquad u,v,w\in\Gamma(-,\gamma\mathcal{M}). \tag{8}
\]
La loi $[\ ,\ ]$ sur $E$ est transportée de la loi $[\ ,\ ]$ sur $\gamma\mathcal{M}$, i.e.\ est définie par la loi
\[
2\,x\wedge y\wedge z=\varphi([x,y],z)\,\omega\qquad \bigl(x,y,z\in\Gamma(-,E)\bigr) \tag{8}
\]
\note{le numéro (8) est porté deux fois sur la page.}
\marginal{\ill{} cf.\ remarque \ill{} que \ill{} précédente}
où $\varphi$ est la forme bil.\ sym.\ associée à $q$, i.e.\ définie par
\[
q(x+y)=q(x)+q(y)+\varphi(x,y). \tag{9}
\]

Nous nous intéressons à la \emph{quadrique affine}
\[
\ \Sigma_q\subset\check{\mathbb{V}}(E)\quad\text{d'équation } q(x)=-1 \tag{10}
\]
\note{un symbole biffé et illisible dans cette formule est omis.}
\note{le $-$ de $-1$ est ajouté à l'encre bleue, comme dans (11).}
qui \struck{correspond à la quadrique} \add{est \uncertain{canon.} iso.\ à la quadrique affine}
\[
\Sigma_{\mathcal{M}}\subset\check{\mathbb{V}}(\underbrace{\gamma\mathcal{M}}_{\mathcal{V}})\quad\text{d'équation } q_{\mathcal{M}}(u)=-1\quad\text{i.e.}\quad \det u=1 \tag{11}
\]
\marginal{$\Sigma_{\mathcal{M}}$ est formé des $u$ \struck{\ill{}} \add{tels que} \struck{$u^2+1=0$} \ill{} $[u^2+1=0]$ des H.C.\ de $A$ \ill{}}
\struck{\ill{}} \uncertain{à} \ill{} \uncertain{correspond} \ill{} affine
\[
\operatorname{Sym}^{*}(\check{E})/(q+1)\operatorname{Sym}^{*}(\check{E})\simeq\operatorname{Sym}^{*}\bigl((\gamma\mathcal{M})^{\vee}\bigr)/(q_{\mathcal{M}}+1) \tag{12}
\]
\note{sous $\gamma\mathcal{M}$ : « $\mathcal{V}\simeq\mathcal{M}/\underline{\mathcal{O}}_S$ ».}

Nous voulons \uncertain{mettre} la \ill{} en relation \ill{} la droite projective définie par la quadrique $q(x)=0$, soit $X$, \struck{\ill{}} \uncertain{choix} \uncertain{de} \uncertain{la} \uncertain{donnée},

\page{91}
\uncertain{ou, \ill{}}, équivalent : la donnée de $(E,q,\omega)$, \ill{} \ill{} : celle de $\mathcal{M}$. Introduisons
\[
S[\underline{i}]=\operatorname{Spec}\underbrace{\underline{\mathcal{O}}_S[\underline{i}]}_{=\underline{\mathcal{O}}_S[T]/(T^2+1)}, \tag{13}
\]
\note{un symbole biffé entre $S$ et $=\operatorname{Spec}$.}
On va définir une immersion ouverte dominante
\[
\Sigma\hookrightarrow X^{S[\underline{i}]}=\underline{\operatorname{Hom}}_S(S[\underline{i}],X) \tag{14}
\]
Pour cela, \struck{explicitons les sections de $X^{S[i]}$ sur $S$} \add{\uncertain{notations}}

\struck{\ill{} \uncertain{préliminaires}}
\[
\begin{cases}
X\simeq\widetilde{X}/\mathbb{G}_m \\
\widetilde{X}\subset\check{\mathbb{V}}(E)^{*}\qquad \widetilde{X}\ \text{d'équation } q(x)=0
\end{cases} \tag{15}
\]
(NB $\widetilde{X}$ est un torseur sur $X$ de groupe $\mathbb{G}_{m,X}$, c'est le torseur associé au Module $\underline{\Omega}^1_{X/S}$ de degré \uncertain{$2$}\,…)

Dire d'\ill{} \ill{} \ill{} que $\widetilde{X}^{S[\underline{i}]}$ \struck{\ill{}} est un $\mathbb{G}_m^{S[\underline{i}]}$-torseur sur $X^{S[\underline{i}]}$, \ill{} les sections \add{de $\widetilde{X}^{S[\underline{i}]}$} sur $S$ s'explicitent comme des sections de $\check{\mathbb{V}}(E)^{*}_{S[\underline{i}]}$ i.e.\ \add{$\mathbb{V}(E)_{S[\underline{i}]}$} \struck{des couples}
\[
\xi=x+\underline{i}y\qquad x,y\in\Gamma(S,E),
\]
et il faut écrire
a) $x+\underline{i}y$ \struck{une section de $E_{S[\underline{i}]}$ \ill{} \ill{} \ill{}} \add{\uncertain{non nulle}} \ill{} fibre ($\Rightarrow x$ et $y$ \ill{}
\note{la fin de la ligne a) et la ligne (16) sont barrées ; on y lit « \ill{} simultanément \ill{} \ill{} fibre ».}

\page{92}
b) $q(x+\underline{i}y)=0$\quad i.e.\quad $q(x)-q(y)+\underline{i}\,\varphi(x,y)=0$
i.e.
\[
\begin{cases} q(x)=q(y) \\ \varphi(x,y)=0 \end{cases} \tag{17}
\]
Il faut faire agir là-dessus $\mathbb{G}_m^{S[\underline{i}]}$, qui s'identifie aux $a+\underline{i}b$ tels que $\underbrace{a^2+b^2}_{\operatorname{Nm}(a+\underline{i}b)}\in\Gamma(S,\mathcal{O}_S^{*})$, opérant sur $x+\underline{i}y$ par
\[
\underbrace{(a+\underline{i}b)}_{\alpha}\underbrace{(x+\underline{i}y)}_{\xi}=(ax-by)+\underline{i}(bx+ay)
\]
Considérons l'ouvert $\widetilde{\mathcal{U}}$ de $\widetilde{X}^{S[\underline{i}]}$ défini par
\[
\ q(x)=q(y)\in\Gamma(S,\mathbb{G}_m) \tag{18}
\]
\note{un symbole biffé et illisible dans cette formule est omis.}
On a, si $\xi=x+\underline{i}y\in\Gamma\widetilde{X}^{S[\underline{i}]}$ et $\alpha=a+\underline{i}b\in\Gamma\mathbb{G}_m^{S[\underline{i}]}$
\[
q(\alpha\xi)=q(ax-by)=a^2q(x)+b^2q(y)-2ab\underbrace{\varphi(x,y)}_{0}=(a^2+b^2)\,q(x)\in\Gamma\mathbb{G}_m
\]
donc $\widetilde{\mathcal{U}}$ est stable par $\mathbb{G}_m^{S[\underline{i}]}$, et définit donc \struck{\ill{}} par passage au quotient un ouvert
\[
\mathcal{U}=\widetilde{\mathcal{U}}/\mathbb{G}_m^{S[\underline{i}]}\subset X^{S[\underline{i}]} \tag{19}
\]
Considérons
\[
\widetilde{\mathcal{U}}\longrightarrow\ \check{\mathbb{V}}(E)
\]
\note{un symbole biffé et illisible dans cette formule est omis.}
défini par
\[
\xi=x+\underline{i}y\overset{\tilde f}{\longmapsto}\ q(x)^{-1}[x,y] \tag{20}
\]
\note{un symbole biffé et illisible dans cette formule est omis.}

\page{93}
Notons que si $x,y$ sont des sections de $E$, on a
\[
q([x,y])=-4q(x)q(y)\underbrace{\sin^2(\widehat{x,y})}_{\overset{\text{déf}}{=}\ \bigl(1-\frac{1}{4}\varphi(x,y)^2/q(x)q(y)\bigr)} \tag{21}
\]
\note{un symbole biffé et illisible dans cette formule est omis.}
\marginal{attention au signe et au facteur 4 !}
\note{dans (21) le $4$ est écrit sur un $1$, et le $\frac14$ sur un autre chiffre.}
i.e.
\[
\text{\struck{$4$}}\ q([x,y])=-4q(x)q(y)+\varphi(x,y)^2 \tag{21 bis}
\]
\[
\text{\struck{$2\varphi([x,y],[x,y])=4(x\wedge y\wedge[x,y])$}}
\]
\note{ligne biffée en zigzag, lecture approximative.}
ce qui, puisque $\varphi(x,y)=0$, \add{$q(x)=q(y)$}, donne
\[
\ q([x,y])=-4q(x)^2
\]
\note{un symbole biffé et illisible dans cette formule est omis.}
ou encore, en car.\ \add{résid.} $\neq 2$,
\[
q\Bigl(\frac{\text{\struck{$4$}}\,[x,y]}{2q(x)}\Bigr)=-1
\]
en d'autres termes \ill{}
\[
\tilde f:\widetilde{\mathcal{U}}\longrightarrow\Sigma_{1} \tag{22}
\]
\note{l'indice $1$ de $\Sigma$ dans (22) est une lecture douteuse.}
\[
\tilde f(\underbrace{x+\underline{i}y}_{\xi})=[x,y]/2q(x)
\]
Je dis que l'on a, pour \add{$\alpha=a+\underline{i}b$} $\alpha\in\Gamma\mathbb{G}_m^{S[\underline{i}]}$
\[
\tilde f(\alpha\xi)=\tilde f(\xi).
\]
En effet \struck{\ill{}}
\[
\alpha\xi=(ax-by)+\underline{i}(bx+ay)=x'+\underline{i}y'
\]
d'où

\page{94}
\[
[x',y']=[ax-by,\,bx+ay]=(a^2+b^2)[x,y]
\]
\struck{\ill{}}
\[
q(x')=q(y')=q(ax-by)=(a^2+b^2)\,q(x)
\]
donc
\[
[x',y']/2q(x')=[x,y]/2q(x).
\]
Par suite (22) se factorise en un morphisme
\[
\mathcal{U}\xrightarrow{\ f\ }\Sigma,\qquad \mathcal{U}\subset X^{S[\underline{i}]} \tag{23}
\]
(si $2$ inversible sur $X$),

\textbf{Proposition} (Le morphisme précédent est un isomorphisme.

a) C'est un monomorphisme. En effet, supposons $\tilde f(\xi)=\tilde f(\xi')$, où $\xi=x+\underline{i}y$, $\xi'=x'+\underline{i}y'$, i.e.
\[
[x,y]/2q(x)=[x',y']/2q(x'),
\]
prouvons que $\xi'$ est de la forme $\alpha\xi$, $\alpha\in\Gamma\mathbb{G}_m^{S[\underline{i}]}$. Notons que $x$ et $y$ ne \add{\uncertain{sont} \uncertain{pas} \uncertain{colin.}} sont pas coll.\ sur aucune fibre (car si on avait $y=\lambda x$ on aurait $\varphi(x,y)=\lambda\varphi(x,x)=2\lambda q(x)=0$ donc $q(x)=0$, \uncertain{absurde}).

\page{95}
D'autre part $x,y$ sont orthogonaux : $x\perp y$ donc : $t=\tilde f(\xi)$, qui n'est pas orth.\ à lui-même en aucune fibre, car $\varphi(t,t)=2q(t)=-2\neq 0$ en chaque fibre. Donc $x,y$ forment une base de $t^{\perp}$, de même $(x',y')$. On doit donc avoir
\[
\begin{cases} x'=ax-by \\ y'=cx+dy \end{cases}\qquad\text{avec } ad-bc\in\Gamma\mathbb{G}_m
\]
\note{le signe de $by$ dans la première ligne est surchargé.}

\struck{\ill{} \ill{}} Écrivons alors $q(x')=q(y')$, $\varphi(x',y')=0$, \ill{} \ill{}
\[
\varphi(x',y')=ac\,\varphi(x,x)-bd\,\varphi(y,y)=0\quad\text{i.e. comme}\quad \varphi(x,x)=2q(x)=\varphi(y,y)=2q(y)
\]
nécessairement
\[
ac-bd=0\qquad\qquad \begin{matrix} a & b\\ d & c\end{matrix}
\]
donc $c=\lambda b$, $d=\lambda a$, avec $\lambda\in\Gamma\mathbb{G}_m$

\struck{et \ill{}} d'où
\[
\begin{cases}
q(x')=(a^2+b^2)\,q(x)\\
q(y')=(c^2+d^2)\,q(y)=\lambda^2(a^2+b^2)\,q(x)
\end{cases}
\]
donc la condition $q(x')=q(y')$ $\in\Gamma(\mathbb{G}_m)$ signifie que
\[
\text{\struck{$(\lambda^2-1)(a^2+b^2)=0$}}
\]
\[
\begin{cases} a^2+b^2\in\Gamma\mathbb{G}_m \\ \lambda^2=1 \end{cases}
\]
\struck{donc} \struck{$(x',y')=$} \struck{$\xi'=x'+\underline{i}y'=\lambda$}

\page{96}
Écrivons enfin
\[
x'\wedge y'/2q(x')=x\wedge y/2q(x)
\]
On a
\[
\begin{cases}
x'\wedge y'=(ad+bc)\,x\wedge y=\lambda(a^2+b^2)\,x\wedge y\\
q(x')=(a^2+b^2)\,q(x)
\end{cases}
\]
donc
\[
x'\wedge y'/q(x')=\lambda\,x\wedge y/q(x)
\]
donc la condition est $\lambda=1$. Donc
\[
\xi'=\alpha\xi,\quad \alpha=a+\underline{i}b,\quad a^2+b^2\in\Gamma\mathbb{G}_m,
\]
donc \emph{mono}.

b) C'est un \emph{épi}, en d'autres termes \struck{localement pour} $\tilde f$ est un épi, ou encore \uncertain{$\forall$} $t\in\Gamma\Sigma$ est loc.\ de la forme $\tilde f(\xi)$. Or on a $\varphi(t,t)=2q(t)\in\Gamma\mathbb{G}_m$, donc $t^{\perp}$ définit un sous-fibré \add{$F$} de $E$ de rang $2$, \struck{donc} supplémentaire de $\mathcal{O}_S\cdot t$, \struck{tout} et la forme induite par $q$ y est non dég. Comme on est en car.\ $\neq 2$, loc.\ pour la top.\ étale on trouve une base orthonormale $x,y$ de $F$.
\[
q(x)=q(y)=1,\quad \varphi(x,y)=0
\]

\page{97}
Considérons alors
\[
t'=\tilde f(x+\underline{i}y)=x\wedge y/q(x)=x\wedge y
\]
\add{et prenons \ill{} \ill{}} il est \uncertain{orth.}\ à $F$ (donc $t$ de la forme $t=\lambda t'$, de plus $q(t)=q(\lambda t')=-1$, donc $\lambda^2=1$. \struck{Considérons}
\[
\text{Donc}\quad
\begin{cases}
q(x)=q(\lambda y)=1,\quad \varphi(x,\lambda y)=0\\
\text{\struck{$x+\underline{i}\lambda y$}}\ \tilde f(x+\underline{i}\lambda y)=x\wedge\lambda y=\lambda\,x\wedge y=\lambda t'=t
\end{cases}
\qquad \text{OK.}
\]
\note{sous la seconde ligne, un début biffé « $f(x+\underline{i}\lambda y)=$ ».}

\textbf{Corollaire}. En inversant l'iso $f$, on trouve une immersion ouverte (en car.\ résid.\ $\neq 2$)
\[
\Sigma\overset{g}{\hookrightarrow}X^{S[\underline{i}]} \tag{24}
\]

\struck{\textbf{Remarque}. Soit $q'=-q$, \ill{}}
\note{cette remarque, jusqu'au bas de la page, est barrée de deux longs traits obliques.}
\[
\text{\struck{$\delta(q')=-\delta(q)=\omega^{\otimes 2}$}}
\]
\struck{donc $(E,q',\omega)$ est} \add{\struck{structure}} \struck{anti} \struck{spéciale orthogonale \uncertain{sans} modification, \ill{} \ill{} \ill{}}
\[
\text{\struck{$\varphi'(x,y)=-\varphi(x,y)$}}
\]
\struck{Si on désigne de façon \ill{} $[\ ,\ ]=[\ ,\ ]_q$ \ill{} par $q'$, on trouve $[\ ,\ ]_{q'}=[\ ,\ ]'$ défini par}
\[
\text{\struck{$\varphi'([x,y]',z)\,\omega=2\,x\wedge y\wedge z$}}
\]
\[
\text{\struck{$\varphi([x,y],z)\,\omega=-\varphi'([x,y],z)\,\omega$}}
\]
\struck{i.e.}
\[
\text{\struck{$[x,y]'=-[x,y]$}}
\]
La formule (21) devient alors
\[
q'([x,y]')=4q'(x)q'(y)\underbrace{\sin^2(\widehat{x,y})'}_{1-\varphi'(x,y)^2/4q'(x)q'(y)} \tag{25}
\]

\page{98}
Notons qu'on a une immersion canonique
\[
X\hookrightarrow X^{S[\underline{i}]}
\]
qui ici s'explicite à partir de
\[
\widetilde{X}\longrightarrow\widetilde{X}^{S[\underline{i}]},\qquad x\longmapsto x+\underline{i}\,0
\]

\note{le cadre qui suit est barré de traits obliques.}
\struck{\textbf{Lemme}} \struck{Soit $\xi=x+\underline{i}y\in\Gamma\widetilde{X}^{S[\underline{i}]}$} \add{\struck{i.e.\ $q(x)=q(y)$, $\varphi(x,y)=0$, $x+\underline{i}y$ partout non nul}}\struck{, Pour que $\xi\in\Gamma\widetilde{\mathcal{U}}$ i.e.\ $q(x)=q(y)\in\Gamma\mathbb{G}_m$, il f.\ et s.\ que $\exists\,\alpha=a+\underline{i}b\in\Gamma\mathbb{G}_m^{S[\underline{i}]}$ ($a^2+b^2\in\Gamma(\mathbb{G}_m)$) et $x_0\in\Gamma\widetilde{X}$ ($q(x_0)=0$, $x_0$ \ill{} \ill{} \ill{} fibre \ill{}) tels que $\xi=\alpha x_0$ i.e.}
\[
\text{\struck{$x=ax_0,\quad y=bx_0$}}
\]
\note{tel qu'écrit, l'énoncé biffé demande « $\xi\in\Gamma\widetilde{\mathcal{U}}$ » ; la conclusion $\xi=\alpha x_0$ correspond plutôt au complémentaire de $\widetilde{\mathcal{U}}$, comme dans le lemme repris plus bas.}

En fait on trouve que $\mathcal{U}\xrightarrow{\ \sim\ }\Sigma$,

\textbf{Proposition}. L'application inverse de $f$ (déduite par passage au quotient de $\tilde f:\widetilde{\mathcal{U}}\to\Sigma$, $\tilde f(x+\underline{i}y)=[x,y]/q(x)$) induit un isomorphisme
\[
\boxed{\ \Sigma\xrightarrow{\ \sim\ }\widetilde{X}^{S[\underline{i}]}-\widetilde{X}=\mathcal{U}\ }
\]
Ceci équivalant au lemme

\textbf{Lemme}. Si $S=\operatorname{Spec}k$, \add{$k$ \uncertain{alg.\ clos} de car.\ $\neq 2$}, \ill{} \ill{} \ill{} \ill{} \struck{\ill{}} \add{\ill{} \ill{} \ill{}}

$x+\underline{i}y\in\Gamma\widetilde{X}^{S[\underline{i}]}$ [$x,y$ deux él.\ de $E$ \struck{tels que $x,y$ non colin., soit \ill{} \ill{}} \add{i.e.\ $y=\lambda x$,} \add{et $q(x)=q(y)$, $\varphi(x,y)=0$}, \struck{\ill{} \ill{} avec} $\lambda^2\neq-1$] \struck{\ill{}}, on a $x+\underline{i}y\notin\widetilde{\mathcal{U}}^{S[\underline{i}]}$ i.e.\ $q(x)=q(y)=0$ ssi
\[
x+\underline{i}y=(a+\underline{i}b)\,x_0,
\]
avec $x_0\in\Gamma\widetilde{X}$ i.e.\ $x_0\in E^{*}$, $q(x_0)=0$, \struck{\ill{}} $a+\underline{i}b\in\Gamma\mathbb{G}_m^{S[\underline{i}]}$ i.e.\ $a^2+b^2\neq 0$.

\page{99}
En effet, \add{soit qu'\ill{}} on a $y=\lambda x$, i.e.\ $x+\underline{i}y=(1+\underline{i}\lambda)x_0$, avec $x_0=x$ et $a+\underline{i}b=1+\underline{i}\lambda$ et on a $a^2+b^2=1+\lambda^2\neq 0$ puisque l'on doit avoir $\lambda^2\neq -1$, d'où nécessité. Pour suffisance, si $x+\underline{i}y=x_0(a+\underline{i}b)$, i.e.\ $x=ax_0$, $y=by_0$, on a bien \ill{} $q(x)=q(y)=a^2q(x_0)=0$ \ill{}

O.K.\,…

\drawing{dans la marge gauche, un cône (triangle et ellipse) dont le sommet est marqué d'une étoile.}

\note{tout ce qui suit sur la page est barré de grands traits obliques.}
\struck{\emph{Cas de $Q$ car.\ $2$}} \add{\struck{$\underline{i}=1+\varepsilon$ \ill{} \ill{} \ill{} \ill{}}} \struck{$S^{[\underline{i}]}$ est alors \ill{} \ill{} des \ill{}, \ill{} $X^{S[\underline{i}]}$ est le fibré tangent à $X$.}

\struck{On trouve}
\[
\text{\struck{$x+\underline{i}y=\underbrace{(x+y)}_{z}+\varepsilon y\longmapsto[x,y]/q(x)$}}
\]
\[
\text{\struck{$=[x+y,y]/q(y)$}}\qquad \text{\struck{$(x+y\neq 0\ \text{i.e.}\ x\neq y)$}}
\]
\[
\text{\struck{$z+\varepsilon y\longmapsto[z,y]/q(y)$}}
\]
\[
\text{\struck{$z\neq 0,\quad q(y+z)=q(y)\neq 0,\quad q(z)=\varphi(y,z)$}}
\]

\struck{\ill{}}
\[
\text{\struck{$\det u=1,\ \operatorname{Tr}u=0\quad(\Rightarrow u^2+1=0)$}}
\]
\[
\text{\struck{$x+\underline{i}y:\ q(x)=q(y)\neq 0,\ \varphi(x,y)=0,\quad y\neq\lambda$}}
\]
\note{ce dernier système, à droite, est précédé d'un point d'interrogation et d'un symbole biffé.}
\[
\text{\struck{$?\Leftrightarrow\ u=[v,w]/\det w$}}
\]
\[
\text{\struck{si $\det v=\det w\neq 0$, $\operatorname{Tr}v=\operatorname{Tr}w=0$, $\operatorname{Tr}vw=0$, $v\neq w$}}
\]
\note{une ligne du système, entre $\operatorname{Tr}vw=0$ et $v\neq w$, est biffée et illisible.}
\[
\text{\struck{$\det[v,w]$}}
\]

\page{100}
\note{le passage suivant, jusqu'à (26), est enfermé dans un cadre et semble abandonné.}
Posons $q''=4q'=-4q$, $\omega''=8\omega$, de sorte que
\[
\delta(q'')=-2^6\,\delta(q)=2^6\,\omega^{\otimes 2}=(8\omega)^{\otimes 2}=\omega''^{\otimes 2}
\]
\[
\delta(q'')=\omega''\cdot\omega''
\]
la formule (25) prend la forme
\[
q''(\ \dots \tag{26}
\]
\note{la formule (26) s'interrompt après « $q''($ ».}

\struck{Remarque} Posons
\[
q_\lambda=\lambda q\qquad \varphi_\lambda=\lambda\varphi\qquad [\ ,\ ]_{\lambda\mu}=(\lambda\mu)^{-1}[\ ,\ ],\qquad \omega_\mu=\mu\omega
\]
\struck{La formule (8) s'écrit alors}
\[
\text{\struck{$q_\lambda([x,y]_\lambda)=\lambda$}}
\]
\[
\text{\struck{$2\,x\wedge y\wedge z=\varphi_\lambda([x,y]_\lambda,2)\,\omega$}}
\]
donc on a
\[
\begin{cases}
\delta(q_\lambda)=\lambda^3\delta(q)=-\lambda^3\,\omega^{\otimes 2}=-\lambda^3\mu^{-2}\,\omega_\mu^{\otimes 2}\\
2\,x\wedge y\wedge z=\varphi_\lambda([x,y]_{\lambda,\mu},z)\,\omega_\mu
\end{cases}
\]
\[
\begin{cases}
q_\lambda([x,y]_{\lambda,\mu})=\lambda(\lambda\mu)^{-2}q([x,y])\\
\qquad =-4\lambda^{-1}\mu^{-2}q(x)q(y)\sin^2(\widehat{x,y})\\
\qquad =-4\lambda^{-3}\mu^{-2}q_\lambda(x)q_\lambda(y)\sin^2(\widehat{x,y})
\end{cases}
\]
On veut choisir $\lambda,\mu$ tels que $-4\lambda^{-3}\mu^{-2}=1$
\drawing{dans la marge, un petit parallélogramme de vecteurs, en partie pointillé.}
i.e.\ $\lambda^{-3}\mu^{-2}=-\frac14$ avec $\lambda,\mu$ des \add{(au signe près)} puissances de $2$
\[
\lambda=-2^{\alpha},\ \mu=2^{\beta}\qquad \lambda^{-3}\mu^{-2}=-2^{-(3\alpha+2\beta)}=-2^{2}
\]
\note{tel qu'écrit : « $=-2^{2}$ » ; la condition $-\frac14$ demanderait $-2^{-2}$.}
i.e.\ \ill{} \ill{}
\[
3\alpha+2\beta=-2\qquad \text{p.ex.\ on prend } \alpha=0,\ \beta=-1
\]
\note{le signe de $\beta=-1$ est une tache d'encre ; avec $3\alpha+2\beta=-2$ et $\alpha=0$ on a bien $\beta=-1$.}
ce qui donne
\[
\lambda=-1,\quad \mu=\tfrac12\qquad q_\lambda=q'=-q,\quad \omega_\mu=\omega'=\tfrac12\omega
\]
\note{$\mu=\frac12$ est écrit sur un « $4$ » ; un indice est biffé sous $q'$ ; la lecture de $\omega_\mu$ est incertaine (un chiffre surchargé).}
\[
\begin{cases}
\delta'(q')=+4\,\omega'^{\otimes 2}=(2\omega')^{\otimes 2}\\
x\wedge y\wedge z=\frac12\,\varphi'([x,y]',z)\,\omega'\\
q'([x,y]')=q'(x)q'(y)\sin^2(\widehat{x,y})\\
\qquad =q'(x)q'(x)\bigl(1-\tfrac14\varphi(x,y)^2\ \dots\bigr)
\end{cases}
\]
\[
[x,y]_{\lambda\mu}=[x,y]'=-2[x,y]
\]
\note{la dernière ligne de l'accolade porte $q'(x)q'(x)$ sur la page ; la fin du facteur est une tache. L'argument se poursuit après le lot.}

\end{document}
